\documentclass{leveldoc}
\DocUses{foundations,deflection_plates,injection}
\DocFingerprint{a28f747aa7d73b98}

\begin{document}

\section{What the level asks}\label{sec:asks}

\leveldescription

One positive ion leaves the left edge at an upward angle. The steering charge
of Injection is built in, below a grounded metal wall, and the wall lies
between that charge and the beam. The player adds charges on the beam's side
of the wall. There are no magnets and no power supplies.

The ion must enter the detector moving along $+x$, inside a cone. The shot
has no kinetic window and no gate. A flight ends when the ion touches a
charge or the wall, leaves the world, or enters the detector
(\xsecref{foundations}{sec:board}). Entering the box counts only when the
direction lies inside the cone (\xeqref{injection}{eq:cone}); otherwise the
outcome is Rejected. The wall's obstacle is its box plus a contact shell of
thickness $\val{contact}$. Each charge's obstacle radius is $\val{Rcharge}$.
The ion's radius is $0$.

The reference adds one charge above the wall. The ion climbs less steeply and
meets the bottom edge of the detector inside the cone. With no added charge
the built-in charge is screened, the ray stays near its launch direction, and
the detector rejects it. This document uses the deflection plates' rectangle
collocation. Both charges are external, on the right-hand side. There is one
wall, so there is no uniform gap. The flight is integrated by the classical
Runge--Kutta method.

\section{Setup}\label{sec:setup}

\setuptable

The speed of light is $c = \val{cLight}$. The ion has charge $q = \val{qPart}$,
mass $m = \val{mPart}$ and radius $0$, and it is launched with kinetic energy
$T_0 = \val{T0}$. The launch is at $(\val{xLaunch}, \val{yLaunch})$, at
$\val{launchAng}$ degrees to $+x$. Then $\gamma_0 = \val{gamma0}$,
$\abs{\vect{p}_0} = \val{p0}$ and $v_0 = \val{v0} = \val{vOverC}\,c$
(\xeqref{foundations}{eq:launch}).

The built-in charge is $Q = \val{qBuilt}$ at $(\val{xBuilt}, \val{yBuilt})$.
The wall is centred at $(\val{xPlateC}, \val{yPlateC})$, grounded at
$V = \val{vPlate}$, with length $\val{plateL}$ along $x$, thickness
$\val{plateT}$ and height $\val{plateH}$. It occupies
$\val{xPlate0} \le x \le \val{xPlate1}$,
$\val{yPlate0} \le y \le \val{yPlate1}$. That slab lies between the built-in
charge and the launch. The reference charge is $Q = \val{qSteer}$ at
$(\val{xSteer}, \val{ySteer})$, on the beam's side of the wall. The largest
magnitude on offer is $\val{qMag}$.

The detector is $\val{xDetLo} \le x \le \val{xDetHi}$,
$\val{yDetLo} \le y \le \val{yDetHi}$. Its cone has axis $\val{axisDeg}$
degrees and half-angle $\val{halfDeg}$ degrees. The world is
$-1 \le x \le \val{xWorld}$, $-1 \le y \le \val{yWorld}$. The reported time
step is $h = \val{H}$. The mesh uses rectangles of edge $\val{panelFine}$,
$\val{nPanels}$ of them on the upper half of the wall.

\begin{figure}[ht]
\centering
\begin{tikzpicture}
\begin{axis}[axis equal image, width=0.92\linewidth,
  xmin=-1.5, xmax=32.5, ymin=-2, ymax=22.5,
  xlabel={$x$}, ylabel={$y$}, tick label style={font=\small}]
  \draw[black!40, dashed] (-1,-1) rectangle (\rawval{xWorld},\rawval{yWorld});
  \draw[black!25] (0,0) rectangle (30,20);
  \fill[green!50!black, fill opacity=0.18]
    (\rawval{xDetLo},\rawval{yDetLo}) rectangle (\rawval{xDetHi},\rawval{yDetHi});
  \draw[green!40!black]
    (\rawval{xDetLo},\rawval{yDetLo}) rectangle (\rawval{xDetHi},\rawval{yDetHi});
  \draw[green!50!black]
    (\rawval{coneOx},\rawval{coneOy}) -- (\rawval{coneXp},\rawval{coneYp});
  \draw[green!50!black]
    (\rawval{coneOx},\rawval{coneOy}) -- (\rawval{coneXm},\rawval{coneYm});
  \fill[black!20]
    (\rawval{xPlate0},\rawval{yPlate0}) rectangle (\rawval{xPlate1},\rawval{yPlate1});
  \draw[black!70]
    (\rawval{xPlate0},\rawval{yPlate0}) rectangle (\rawval{xPlate1},\rawval{yPlate1});
  \fill[violet] (\rawval{xBuilt},\rawval{yBuilt}) circle[radius=\rawval{Rcharge}];
  \fill[violet] (\rawval{xSteer},\rawval{ySteer}) circle[radius=\rawval{Rcharge}];
  \addplot[orange!80!black] table {pathbare.dat};
  \addplot[very thick, blue!60!black] table {pathref.dat};
  \draw[-{Stealth[length=5pt]}, thick]
    (\rawval{xLaunch},\rawval{yLaunch}) -- (\rawval{arrowX},\rawval{arrowY});
  \fill (\rawval{xLaunch},\rawval{yLaunch}) circle[radius=0.08];
\end{axis}
\end{tikzpicture}
\caption{The wall, the two charges and the two flights. The gray slab is the
grounded wall. Both discs are negative: the lower one is built in, and the
upper one is the reference charge. The green box is the detector, and the
green wedge is the accepted cone about $+x$. Blue is the reference: the ion
meets the bottom edge. Orange keeps the wall and the built-in charge and
meets the left edge.}
\label{fig:board}
\end{figure}

\section{Principles and model}\label{sec:model}

\begin{principles}
  \item The relativistic equation of motion \xeqref{foundations}{eq:eom}, the
    launch \xeqref{foundations}{eq:launch}, and energy in a static field
    \xeqref{foundations}{eq:energy}. There is no magnetic field.
  \item The wall is a grounded metal box (\physref{2.7}). Its surface charge
    is the deflection plates' collocation
    \xeqref{deflection_plates}{eq:collocation}, with the rectangle integral
    \xeqref{deflection_plates}{eq:rectangle}. Both charges are external.
  \item The detector counts an entry only inside the cone
    \xeqref{injection}{eq:cone}.
  \item The board's rules (\xsecref{foundations}{sec:board}) and the classical
    Runge--Kutta method (\xsecref{foundations}{sec:numerics}).
  \item The radiated power \xeqref{foundations}{eq:larmor}, summed by Simpson's
    rule \xeqref{foundations}{eq:simpson}. It is not put back into the force.
\end{principles}

The reference has the grounded wall, the built-in charge and the reference
charge. The state obeys \xeqref{foundations}{eq:eom} with $\vect{B} = \vect{0}$
until the first ending event. The cone is applied to that event when it is an
entry.

\begin{assumption}{A static wall}\label{as:static}
The potential equals the grounded value at once, and the charges are fixed.
The launch speed is $\val{vOverC}\,c$. The game uses the same static wall
(\physref{2.7}).
\end{assumption}

\begin{assumption}{No radiation reaction}\label{as:radiation}
The radiated energy is not taken from the flight. On the reference flight it
is a fraction $\val{RefRadFrac}$ of $T_0$.
\end{assumption}

\begin{assumption}{No image of the ion}\label{as:image}
The ion's own surface charge is left out, as on the deflection plates and as
in the game (\physref{2.7}). The bound $q^2/d^2$, with $d$ floored at
$10^{-3}$ and taken relative to $\max(\abs{q\vect{E}}, T_0)$, reaches
$\val{RefImage}$ on the reference samples and $\val{OffImage}$ with no added
charge. Both lie under the engine's limit $10^{-10}$.
\end{assumption}

\begin{assumption}{Equal rectangles}\label{as:panels}
The mesh is equal rectangles of edge $\val{panelFine}$. Replacing the edge by
$\val{panelCoarse}$ moves an end state by $\val{errMesh}$. Both flights use
the wall, so both enter that estimate.
\end{assumption}

\begin{assumption}{The fixed step}\label{as:step}
The reported step is $h = \val{H}$. Halving it estimates an end state by
$\val{errStep}$. The tolerance on a state is ten times the larger of that
estimate and the panel-size estimate, and at least $10^{-9}$. Here the
panel-size estimate is the larger one, so the tolerance is $\val{tolFlight}$.
Between samples the ion moves at most $\val{stepReach}$. The clearance of the
wall is sampled at the ends of the steps and at the ending point. It reaches
$\val{RefDmin}$ on the reference and $\val{OffDmin}$ with no added charge.
\end{assumption}

\section{Derivation}\label{sec:derivation}

\subsection{The wall}\label{sec:wall}

The rectangles, the mirror and the collocation are
\xeqref{deflection_plates}{eq:rectangle} and
\xeqref{deflection_plates}{eq:collocation}. The wall potential on the
right-hand side is $\val{vPlate}$, and each charge's potential is subtracted
there. The ion is not in the linear system. The reported mesh has
$\val{nPanels}$ upper rectangles. The residual at the centroids is
$\val{solveRes}$. On the long faces, off the centroids, the potential differs
from the grounded value by at most $\val{surfRes}$.

There is no second plate, so the uniform-gap field does not apply. One cell
above the wall, over the built-in charge, the point-charge field is
$\val{eCoulomb}$ along $-y$ and the potential is $\val{phiCoulomb}$. With the
wall, and without the reference charge, the field at the same station is
$\val{eScreen}$ along $+y$ and the potential is $\val{phiScreen}$. The ratio
of the magnitudes is $\val{screenRel}$. The induced charge on the wall and its
mirror is $\val{qWall}$, the fraction $\val{wallFrac}$ of $\abs{Q}$ and of the
opposite sign. On the straight launch line at the same abscissa the
point-charge field is $\val{eRayCoulomb}$ along $-y$, and the wall leaves
$\val{eRay}$, still along $-y$, the fraction $\val{rayRel}$. The march uses
the panels: the ion's clearance of the metal is $\val{RefDmin}$.

\subsection{The cone}\label{sec:cone}

The test is \xeqref{injection}{eq:cone}. The axis is $+x$ and the half-angle
is $\val{halfDeg}$ degrees. The momentum at the geometric entry gives
$\beta$. The entry counts when $\beta$ is at most that half-angle. A larger
$\beta$ rejects the ion at the entry point, and the flight stops there.

\subsection{The flight}\label{sec:flight}

The diagnostic is \xeqref{deflection_plates}{eq:flight-energy}. The state is
advanced by \xeqref{foundations}{eq:rk4} with step $h = \val{H}$. Events are
located as on the deflection plates: thirty-two Hermite samples, then sixty
bisections, obstacles before the world, and the world before the detector.
The cone is applied after that geometric entry. The step-halving estimate
divides the change by $15$; the panel-size estimate is the change itself when
the edge doubles.

On the straight line above the reference charge, that charge's field is
$\val{eSteer}$ along $-y$, and the field of the wall plus the built-in charge
is $\val{ePipe}$. The ratio is $\val{steerRel}$. A positive ion is pulled
downward by the reference charge, and that field dominates the screened
charge.

\section{The reference solution by hand}\label{sec:hand}

\begin{procedure}
  \item\label{st:launch} \emph{The launch} \xeqref{foundations}{eq:launch}.
    $\gamma_0 = \val{gamma0}$, $\abs{\vect{p}_0} = \val{p0}$ and
    $v_0 = \val{v0}$. The direction is the file's, at $\val{launchAng}$
    degrees to $+x$.
  \item\label{st:wall} \emph{The wall.} Keep $Q = \val{qBuilt}$ at
    $(\val{xBuilt}, \val{yBuilt})$. The wall is centred at
    $(\val{xPlateC}, \val{yPlateC})$, at potential $\val{vPlate}$. Solve
    \xeqref{deflection_plates}{eq:collocation}. The centroid residual is
    $\val{solveRes}$ and the off-centroid discrepancy is $\val{surfRes}$.
    The induced charge is $\val{qWall}$.
  \item\label{st:steer} \emph{The reference charge.} Place
    $Q = \val{qSteer}$ at $(\val{xSteer}, \val{ySteer})$ and solve the
    collocation again. The wall then carries $\val{qWallRef}$.
  \item\label{st:march} \emph{The march.} Advance
    \xeqref{deflection_plates}{eq:flight-energy} and
    \xeqref{foundations}{eq:eom} by \xeqref{foundations}{eq:rk4} with step
    $\val{H}$. The state every five time units:
    \begin{center}\small\val{states}\end{center}
  \item\label{st:end} \emph{The arrival.} The ion meets the \val{RefFace}
    face of the detector at $(\val{RefX}, \val{RefY})$, at $t = \val{RefT}$,
    with momentum $(\val{RefPx}, \val{RefPy})$, at an angle $\val{RefAng}$
    degrees from $+x$. The ordinate is the detector's bottom edge. The cone
    angle is $\val{RefBeta}$ degrees, inside the half-angle $\val{halfDeg}$
    degrees, so the outcome is \val{RefOutcome}. The run takes
    $\val{RefSteps}$ steps. The path only climbs: its lowest point is the
    launch, and its highest is $\val{RefYmax}$.
  \item\label{st:check} \emph{The checks.} $W$ drifts by at most
    $\val{RefWdev}$. The step-halving estimate is $\val{RefErr}$, and the
    larger of the two flights is $\val{errStep}$. The panel-size estimate is
    $\val{errMesh}$. The sampled clearance is $\val{RefDmin}$, outside the
    shell $\val{contact}$. The image-force ratio reaches $\val{RefImage}$.
  \item\label{st:rad} \emph{The radiation.} Simpson's rule on steps of $h/2$
    gives $\val{RefRad}$, a fraction $\val{RefRadFrac}$ of $T_0$. Halving
    that step again changes the integral by the relative amount
    $\val{RefRadErr}$. This integral is not a row of the comparison.
\end{procedure}

\begin{result}[The reference flight]
The ion enters the detector through its bottom edge at
$(\val{RefX}, \val{RefY})$ at $t = \val{RefT}$, at $\val{RefBeta}$ degrees
to $+x$, inside the cone. The level is solved.
\end{result}

\section{Comparison with the game}\label{sec:comparison}

\comparisontable

The table holds the outcome, the time and the end position of the reference,
and the same three for the empty placement. The tolerance on a state is
$\val{tolFlight}$. Ten times the panel-size estimate $\val{errMesh}$ lies
above $10^{-9}$, and it is larger than the step-halving estimate
$\val{errStep}$, so the tolerance is that product. The tolerance on a state
was fixed from these estimates before the comparison.

The largest absolute difference is the reference abscissa,
$\val{cmpRefXDiff}$. It lies above the panel-size estimate and under the
tolerance. The reference time differs by $\val{cmpRefTDiff}$, also above the
panel-size estimate and under the tolerance. The empty-placement time differs
by $\val{cmpOffTDiff}$ and its ordinate by $\val{cmpOffYDiff}$, both under
the panel-size estimate. The outcomes agree.

The reference ordinate agrees exactly. It is the detector's bottom edge. The
empty placement's abscissa agrees exactly. It is the detector's left edge.

\section{Other solutions}\label{sec:others}

The empty placement in \code{placements.json} keeps the wall and the built-in
charge and adds nothing. It does not solve the level. The ion meets the
\val{OffFace} face of the detector at $(\val{OffX}, \val{OffY})$, at
$t = \val{OffT}$, with momentum $(\val{OffPx}, \val{OffPy})$. The cone angle
is $\val{OffBeta}$ degrees, outside the half-angle, so the outcome is
\val{OffOutcome}. The ordinate lies $\val{straightGap}$ above the
straight-ray ordinate $\val{straightY}$. The energy drifts by $\val{OffWdev}$,
and the run takes $\val{OffSteps}$ steps. The sampled clearance is
$\val{OffDmin}$.

The reference charge is what brings the cone angle from $\val{OffBeta}$
degrees down to $\val{RefBeta}$ degrees, and the meeting point from the left
face down to the bottom edge. This document integrates that charge and the
empty placement.

\section{What the game computes}\label{sec:engine}

The wall is a grounded electrode (\physref{2.7}). The surface charge is the
engine's triangle collocation, the same boundary condition as
\xeqref{deflection_plates}{eq:collocation}, with both charges on the
right-hand side. This document uses the equal rectangles of the deflection
plates. The end states move by $\val{errMesh}$ when that edge doubles, and
the reference abscissa in \S\ref{sec:comparison} lies above that estimate and
under ten times it.

The direction test is \xeqref{injection}{eq:cone}. The ion's image force is
omitted on both sides. The bound sampled on the reference flight reaches
$\val{RefImage}$, under the engine's limit $10^{-10}$. The flight in the
engine is the adaptive method of \physref{5} and \physref{6}, verified as in
\physref{7}. Our step estimate is $\val{errStep}$. Radiation is integrated
and not applied to the momentum (\xsecref{foundations}{sec:radiation},
\physref{8}).

\section{Further reading}\label{sec:reading}

The rectangle integral, the collocation and the flight energy are derived for
the deflection plates (\xeqref{deflection_plates}{eq:rectangle},
\xeqref{deflection_plates}{eq:collocation},
\xeqref{deflection_plates}{eq:flight-energy}). The cone is derived for
Injection (\xeqref{injection}{eq:cone}). The motion and the Runge--Kutta step
are in the Foundations (\xsecref{foundations}{sec:numerics}). The engine's
electrodes are \physref{2.7}.

\levelfileappendix

\end{document}
